\chapter{Introducción}
\label{chap:introduccion}

% Versión 1 del capítulo 1 (2026-10-04). Base: plan de trabajo aprobado
% (introducción, planteamiento del problema, justificación y objetivos), redactado
% como informe y ajustado al alcance que sustentan los capítulos 3 a 5.

Conocer la distribución de densidad en el interior de un edificio volcánico es una de las bases de la evaluación de su amenaza. Los conductos por los que asciende el magma, las cavidades, las zonas de alteración hidrotermal y los cuerpos de roca más densa modifican la masa que hay a lo largo de cualquier trayecto que atraviese la montaña, y su localización ayuda a interpretar procesos como la desgasificación, el ascenso de material o la estabilidad de los flancos \parencite{tanaka2007asama,tanaka2023muography}. Los métodos geofísicos convencionales estiman esa estructura de forma indirecta, a partir de campos medidos en la superficie, y las técnicas de inspección por radiación artificial no alcanzan a atravesar cientos de metros de roca. Este trabajo se ocupa de una técnica que sí lo hace, la muografía por transmisión, y de una de sus limitaciones más severas: el tiempo de exposición que exige obtener una imagen interpretable.

\section{Contexto y planteamiento del problema}
\label{sec:intro-problema}

La muografía aprovecha los muones que se producen cuando los rayos cósmicos interactúan con la atmósfera. Por su energía elevada y su interacción débil con la materia, una fracción de esos muones atraviesa cientos de metros o incluso kilómetros de roca, y la atenuación de su flujo en cada dirección de observación indica cuánta materia hay a lo largo de ella \parencite{bonomi2020,tanaka2023muography}. Desde sus primeras aplicaciones volcánicas, la técnica ha permitido visualizar estructuras bajo el cráter del monte Asama \parencite{tanaka2007asama} y obtener imágenes de alta definición del volcán Sakurajima \parencite{olah2018sakurajima}, ambos en Japón. En Colombia se diseñó y construyó el telescopio híbrido MuTe con el propósito específico de estudiar volcanes activos del país en zonas de difícil acceso \parencite{penarodriguez2020}.

La fuente de la técnica es, al mismo tiempo, su principal limitación. A diferencia de la radiografía médica o industrial, en la que la intensidad de la fuente puede ajustarse, la muografía depende del flujo natural de muones, cercano a un muon por centímetro cuadrado y por minuto sobre un detector horizontal a nivel del mar \parencite{pdg2022}. En las direcciones casi horizontales con las que se observa un volcán ese flujo es mucho menor, y después de atravesar cientos de metros de roca solo sobrevive una fracción pequeña de él. El número de muones registrado en cada dirección es, por lo tanto, reducido y está sujeto a fluctuaciones aleatorias, de modo que se requieren exposiciones de semanas o meses para que los contrastes de densidad se distingan de esas fluctuaciones \parencite{lesparre2010}. Durante ese tiempo el instrumento debe operar en ambientes remotos, expuesto a la intemperie y con suministro de energía limitado, lo que encarece la adquisición, aumenta el riesgo de fallas y limita la frecuencia con la que puede actualizarse una imagen \parencite{bonomi2020,lechmann2021}. Las imágenes obtenidas con exposiciones cortas presentan un ruido de conteo elevado que puede enmascarar precisamente las estructuras que se pretende detectar.

En otros campos de la imagenología, el aprendizaje profundo ha ofrecido una alternativa al aumento del tiempo de adquisición. En la tomografía computarizada de baja dosis, por ejemplo, las redes generativas antagónicas se han usado para transformar imágenes ruidosas en estimaciones cercanas a las de dosis completa \parencite{yang2018}. En muografía, \textcite{odonnell2025particles} aplicaron una red generativa antagónica condicional de Wasserstein con penalización de gradiente (cWGAN-GP) a imágenes de baja estadística de bloques de concreto y reportaron que imágenes de un día de exposición alcanzaron, según el índice de similitud estructural, la calidad perceptual de imágenes de 21 días. Ese resultado se obtuvo, sin embargo, en un escenario de infraestructura civil, cuyas geometrías, distancias y contrastes de densidad difieren sustancialmente de los de un edificio volcánico. Otros trabajos recientes han aplicado el aprendizaje automático a la supresión de fondo en muografía volcánica, a la reconstrucción y a la tomografía por dispersión, como se revisa en la \cref{sec:mt-estado-arte}.

De esta situación surgió la pregunta que orientó el trabajo: si un modelo de aprendizaje profundo, entrenado con datos sintéticos generados mediante simulación Monte Carlo, puede transformar imágenes muográficas por transmisión de baja estadística en imágenes de calidad comparable a las obtenidas con exposiciones prolongadas, en escenarios volcánicos y sin perder los contrastes de densidad que dan sentido a la imagen. La respuesta que permiten los resultados, con sus condiciones y sus límites, se presenta en el \cref{chap:conclusiones}.

\section{Justificación}
\label{sec:intro-justificacion}

El trabajo se justificó por la confluencia de una necesidad práctica, una oportunidad metodológica y la disponibilidad de herramientas para abordarla.

La necesidad práctica es el costo del tiempo de exposición. Como se expone en la \cref{sec:mt-estadistica}, la relación entre la señal y el ruido de un conteo crece solo con la raíz cuadrada del número de eventos registrados, de manera que reducir a la mitad el ruido relativo de una imagen exige cuadruplicar su exposición. Cualquier procedimiento que extraiga más información de una adquisición corta incide directamente sobre el costo logístico y sobre la frecuencia con la que puede vigilarse un volcán, y por eso la reducción del tiempo de adquisición es una preocupación reconocida en la comunidad de la muografía \parencite{tanaka2023muography,lechmann2021}.

La oportunidad metodológica es la brecha que deja la literatura. El resultado de \textcite{odonnell2025particles} mostró que la mejora predictiva de imágenes muográficas de baja estadística es factible en estructuras de concreto, y otros trabajos han aplicado redes neuronales a la reducción del ruido estadístico de imágenes muográficas y a la mejora de tomografías \parencite{lefevre2024} o a la tomografía por dispersión de corta duración \parencite{wang2026}. En la literatura revisada no se identificó, en cambio, la aplicación de redes generativas condicionales a imágenes de baja estadística de muografía volcánica por transmisión, cuya geometría de observación y cuyos contrastes difieren de los escenarios estudiados.

La disponibilidad de herramientas hizo viable el estudio. Los pares formados por una imagen de baja estadística y su referencia de alta estadística, necesarios para el aprendizaje supervisado, no existen en la práctica: una medición de campo no tiene una versión equivalente sin ruido, y repetirla con una exposición mucho mayor no garantiza que el volcán permanezca igual. La simulación Monte Carlo del transporte de muones con Geant4 \parencite{agostinelli2003,allison2016}, alimentada por la parametrización del flujo cósmico y el generador EcoMug \parencite{guan2015muonflux,pagano2021}, permite en cambio construir esos pares en condiciones controladas, con la estructura interna de cada escenario conocida de antemano y sin depender de campañas experimentales costosas. Esto hizo posible abordar el problema en el marco de un trabajo de grado y facilita la reproducción de sus resultados. El proyecto se inscribió, además, en la línea de trabajo del Grupo de Óptica y Tratamiento de Señales de la Universidad Industrial de Santander, que articula la simulación física, el procesamiento de imágenes y el aprendizaje profundo, y en el interés nacional por la muografía de volcanes colombianos \parencite{penarodriguez2020}.

\section{Objetivos}
\label{sec:intro-objetivos}

\subsection{Objetivo general}
\label{subsec:intro-objetivo-general}

Desarrollar y validar un esquema de reconstrucción muográfica por transmisión asistido por aprendizaje profundo para reducir el tiempo efectivo de adquisición en exploración volcánica, mediante mejora predictiva de imágenes de baja estadística hacia imágenes de referencia de alta estadística, evaluando la mejora en calidad y detectabilidad de estructuras internas.

\subsection{Objetivos específicos}
\label{subsec:intro-objetivos-especificos}

\begin{enumerate}[label=\textbf{OE\arabic*.}, leftmargin=*, labelsep=0.6em]
  \item Generar un conjunto de datos sintéticos de imágenes muográficas por transmisión, con diferentes niveles de estadística de conteo, mediante simulación Monte Carlo que reproduzca escenarios representativos de exploración volcánica.
  \item Diseñar e implementar un modelo de aprendizaje profundo basado en redes generativas adversarias condicionales para transformar imágenes muográficas de baja estadística en imágenes de calidad comparable a las de alta estadística.
  \item Evaluar el desempeño del modelo entrenado y determinar el factor de reducción del tiempo de adquisición alcanzable con el esquema propuesto, comparando la calidad de las imágenes mejoradas por el modelo frente a imágenes obtenidas con tiempos de exposición prolongados.
\end{enumerate}

Los tres objetivos se abordaron en secuencia, porque el producto de cada uno es el insumo del siguiente. El OE1 se desarrolló mediante una campaña de simulación con Geant4 sobre diecinueve escenarios volcánicos idealizados, de la que se obtuvieron 9\,120 pares de entrada y referencia con cuatro fracciones de conteo, particionados por geometría en doce escenarios de entrenamiento, tres de validación, tres de prueba y uno de generalización (\cref{sec:met-datos,sec:res-oe1}). El OE2 se desarrolló con una red generativa antagónica condicional de Wasserstein con penalización de gradiente, entrenada y evaluada con el 5\,\% de los conteos de referencia; sus resultados respaldan el diseño, la implementación y el entrenamiento del modelo, mientras que la calidad comparable a la de alta estadística queda respaldada solo de forma parcial, porque la mejora de la similitud con la referencia no se acompañó de la conservación de las cuentas regionales (\cref{sec:met-modelo,sec:res-oe2,sec:conc-oe2}). El OE3 se desarrolló mediante la evaluación de los modelos con entradas al 1, 5, 10 y 50\,\% de los conteos y un criterio por vista frente a la entrada sin red al 50\,\%; con ese criterio no se demostró un factor de reducción del tiempo de adquisición, y el factor nominal de 20 que corresponde a la fracción de entrenamiento describe solo la reducción de conteos de la entrada (\cref{sec:met-adquisicion,sec:res-factor,sec:conc-oe3}).

\section{Alcance y delimitación}
\label{sec:intro-alcance}

El plan de trabajo delimitó el estudio a geometrías simplificadas, manejables computacionalmente a nivel de pregrado, y a mapas bidimensionales obtenidos a partir de la atenuación del flujo de muones. Dentro de ese marco, el trabajo se planteó como un estudio enteramente computacional, y sus afirmaciones se refieren a las condiciones que se describen a continuación.

Los escenarios volcánicos son conos idealizados de roca homogénea con anomalías internas de geometría simple, del tipo que la muografía volcánica asocia con conductos, cavidades y cuerpos densos. No reproducen la topografía ni la estructura de densidades de un volcán concreto. Cada imagen es un mapa bidimensional de conteos direccionales, ordenado por elevación y acimut desde una posición de observación, que registra los muones primarios que alcanzan el plano de conteo con la energía y la dirección establecidas; no incluye la respuesta de un detector físico, el fondo instrumental ni las partículas secundarias, y su exposición es una escala convencional y no un tiempo calibrado de medición. En consecuencia, el trabajo no aborda la reconstrucción tridimensional, la inversión de densidad ni la topografía real, y no incluye datos experimentales de campo.

Los modelos de aprendizaje profundo se entrenaron con una sola fracción de conteo, el 5\,\% de la referencia, sobre los pares sintéticos del OE1, y se evaluaron después con el 1, el 5, el 10 y el 50\,\%; su desempeño se refiere a referencias idealizadas de un catálogo conocido. Los indicadores de anomalías se calcularon en regiones de ubicación conocida, por lo que el estudio no aborda la detección ni la localización de anomalías desconocidas. Las conclusiones no se extienden por sí solas a mediciones reales, a otros volcanes ni a modelos entrenados con otras fracciones, en los términos precisados en el \cref{chap:conclusiones}.

\section{Estructura del documento}
\label{sec:intro-estructura}

El documento se organiza en cinco capítulos. El \cref{chap:marco-teorico} presenta los fundamentos sobre los que se apoya el trabajo: la física de los muones cósmicos y su interacción con la roca, la formación de la imagen muográfica por transmisión y sus límites, la estadística de conteo, los modelos de aprendizaje profundo para la mejora de imágenes, el estado del arte de su aplicación en muografía, el papel de los datos sintéticos y las métricas de calidad de imagen. El \cref{chap:metodologia} describe el diseño de la investigación, la generación del conjunto de datos sintéticos, el diseño e implementación del modelo y el protocolo de evaluación. El \cref{chap:resultados} presenta y discute los resultados de cada objetivo específico, y el \cref{chap:conclusiones} reúne las conclusiones por objetivo, la conclusión general y las líneas de trabajo futuro.
